\documentclass[11pt]{article}
\usepackage[a4paper, margin=2.0cm]{geometry}
\usepackage[T1]{fontenc}
\usepackage[utf8]{inputenc}
\usepackage{booktabs, array, amsmath, xcolor, enumitem}
\usepackage[colorlinks=true, linkcolor=blue!50!black, urlcolor=blue!50!black]{hyperref}
\setlength{\parskip}{4pt}
\sloppy

% Every number below comes from experiments/openfhe/energy_campaign_tables/
% (generated by experiments/openfhe/make_energy_campaign_tables.py from
% aggregate.py --scenario=energy_campaign).
\newcommand{\tab}[1]{../../experiments/openfhe/energy_campaign_tables/#1}
\input{\tab{numbers.tex}}

\title{Energy and memory campaign, SEAL and OpenFHE\\[2pt]
\large Summary of the checks and of the two campaign sessions}
\author{Boungo Tr\'esor}
\date{\today}

\begin{document}
\maketitle

\noindent\textbf{Purpose.} As authorized on 7 November: stage~1 (checks) then stage~2 (two paired sessions on two different days, on mains, boost off), reported by session, with a ratio between the libraries only for matched-profile cells with equivalent boundaries. The rules were fixed before the sessions (\texttt{experiments/openfhe/campaign/ENERGY\_CAMPAIGN\_RULES.md}); Sections~5 and~6 of that file record the two changes made during the campaign. Gross package energy is the primary measurement, idle-corrected energy the secondary one, memory the observed peak process RSS. Raw data: \texttt{experiments/results/raw/energy\_campaign/}; every table below is generated from the data.

\paragraph{Summary.}
\begin{itemize}[leftmargin=*, itemsep=2pt]
  \item \textbf{Both sessions are complete and valid}: \nEnergyOne\ energy and \nMemOne\ memory invocations in each, boost never on, highest core frequency \freqMaxTwo~MHz.
  \item \textbf{The measurement is reproducible across the two days}: for \acrossTen\ of the \acrossN\ cells the gross energy of day2 is within 10\% of day1 (\acrossFive\ within 5\%; median day2/day1 \acrossMedian); peak RSS differs by at most \rssAcrossMaxPct\%.
  \item \textbf{OpenFHE/SEAL energy ratios} (\nRatioCells\ matched-profile cells): decryption \rDecNoBgvMin--\rDecNoBgvMax\ (BGV: \rBgvMin--\rBgvMax), multiplication \rMulMin--\rMulMax, relinearization \rRelinMin--\rRelinMax, key generation \rKgMin--\rKgMax. The ratio of each cell is nearly the same in both sessions.
  \item \textbf{Energy does not always follow time}: for decryption, multiplication and relinearization the energy ratio is within \latDevMaxPct\% of the latency ratio of the Standard campaign; for key generation it is clearly higher (e.g.\ \kgEnBig\ against \kgLatBig\ at CKKS~16384/1): SEAL's key generation draws less power (median \kgPowSeal~W against \kgPowOfhe~W), and the time ratio at fixed frequency differs from the one with boost.
  \item \textbf{Peak RSS}: OpenFHE uses more memory at small $N$ (up to $\times$\rssRatioMax) and less at CKKS $N\ge 8192$, category~1 (down to $\times$\rssRatioMin).
\end{itemize}

\section{Stage 1: checks (before the campaign)}

\begin{itemize}[leftmargin=*, itemsep=2pt]
  \item \textbf{Guard test} (session g2): an invocation interrupted by unplugging the charger was set aside with its files kept and measured again after mains returned.
  \item \textbf{Window test} (session w1, \nEnergyW\ energy invocations valid): energy per operation with 20~ms blocks within \winDevPct\% of 50~ms blocks; the fixed cost per block $|c|$ at most \winCPct\% of one OpenFHE key-generation block. Both criteria of the decision rule were met: the single-operation block was accepted, and with the largest relative standard deviation of the OpenFHE key-generation runs (\winRsdPct\%), \keygenReps\ repetitions give a 95\% half-width of \winCiPct\%; key generation keeps \keygenReps\ repetitions.
  \item \textbf{Short check} (session c1, \cellsC\ cells, including configurations outside the validation and the OpenFHE-only operations): \nEnergyC\ energy and \nMemC\ memory invocations valid.
\end{itemize}

\section{Stage 2: the two sessions}

\begin{center}\footnotesize\setlength{\tabcolsep}{4pt}\input{\tab{sessions.tex}}\end{center}

\noindent Flagged inv.: \texttt{high\_idle} (own idle reading $>1.15\times$ the session median); no \texttt{frequency}, \texttt{idle\_window} or \texttt{operator} flag. Pairs $>$10\%: cells whose two energy invocations differ by more than 10\% (\texttt{pair\_disagreement}). Flags remove nothing; every summary is given with all valid runs and without the flagged ones.

\paragraph{Deviations from the plan.} All are recorded in the rules file and in \texttt{set\_aside.csv} of each session; no file was deleted.
\begin{enumerate}[leftmargin=*, itemsep=2pt]
  \item \textbf{Session d1 stopped; short OpenFHE key generations timed in groups.} The first campaign session (d1) stopped, as the validator requires, on a zero energy reading: an OpenFHE key generation at CKKS $N=2048$ lasts about 0.8~ms, close to the counter's update interval. The window test had covered key generations of about 18~ms only. Before any result was analysed, an OpenFHE key generation shorter than 10~ms is timed as $K$ complete key generations in one window of about 50~ms (each from an empty context cache), and divided by $K$; this lowers the time per key generation by about 9\% (warm-up, measured before the campaign: 1.72 against 1.89~ms). d1 is kept and not used. The campaign was run as the new sessions day1 and day2.
  \item \textbf{Time-based set-aside after a recorded disturbance} (rules, Section~6). Decided by the start time of a disturbance with a recorded cause, never by the values; the runs were measured again by resuming the session. \emph{day1}: after an unexplained restart at 06:56, a browser process used 100\% of one core; the \asideDistOne\ energy invocations measured from that time were set aside and measured again. \emph{day2}: a power cut of about 7~hours with repeated sleeps; after mains returned, the battery was recharging from 23\% and idle readings were about twice the session median: the \asideDistTwo\ invocations measured from then on were set aside and measured again with the battery full. One further invocation was set aside automatically when mains was lost. A check of the charger, the power profile and background processes now runs before every start or resume.
  \item \textbf{Clock}: during the power cut the system clock jumped by one hour; it was corrected by network time before the repeat. Only the timestamps of a few set-aside entries are affected, not the measurements.
\end{enumerate}
\noindent \flaggedMorningTwo\ of the \flaggedTwo\ flagged invocations of day2 fall between 06:08 and 06:43; no cause was recorded, so they are flagged and kept, not set aside.

\section{OpenFHE and SEAL}

Ratios OpenFHE/SEAL of gross energy per operation, within each session, for the matched-profile cells (\texttt{probe/operations\_v2.csv}) with an equivalent computational boundary. Rotation has no energy measurement. For encryption, addition and the cells requiring separate analysis, both libraries' values are in \texttt{results/final/energy\_campaign\_comparisons.csv}, without a ratio.

\begin{center}\footnotesize\setlength{\tabcolsep}{3.5pt}\input{\tab{compare.tex}}\end{center}

\noindent day2$^\circ$: without the flagged invocations. $^\dagger$: OpenFHE timed as several key generations per window (deviation~1). Peak RSS is the whole process (library, context, keys, ciphertexts), day1; the RSS ratio is nearly the same in day2.

\paragraph{Observations} (limited to these cells, this laptop and a fixed frequency of about 3.27~GHz).
\begin{itemize}[leftmargin=*, itemsep=2pt]
  \item The median ratio is \ratioMedianOne\ (day1) and \ratioMedianTwo\ (day2); OpenFHE uses less energy in \ratioLowerOne\ and \ratioLowerTwo\ of the \nRatioCells\ cells. With idle-corrected energy the median ratio is \netMedianOne\ and \netMedianTwo: the conclusions do not change.
  \item The largest difference between sessions, CKKS 16384/1 relinearization (\relinBigOne\ and \relinBigTwo), comes from one SEAL invocation of day2 flagged \texttt{high\_idle}; without it the day2 ratio is \relinBigTwoClean.
  \item Key generation: OpenFHE uses less energy than SEAL, but its advantage is smaller in energy than in time (\kgEnBig, \kgEnMid, \kgEnSmall\ against latency ratios \kgLatBig, \kgLatMid, \kgLatSmall). Two effects contribute: SEAL draws less power during key generation (\kgPowSeal\ against \kgPowOfhe~W; other operations: \opPowSeal\ and \opPowOfhe~W), and the time ratio at fixed frequency is not the same as with boost. Energy is therefore reported in its own right, not derived from latency.
  \item Memory: the RSS ratio rises above 1 at small $N$, presumably because each library's fixed footprint dominates there, and falls to about 0.5 at CKKS 16384/1.
\end{itemize}

\paragraph{Next steps.} Thesis draft (results chapter with these tables) for 29~November; one latency pilot each for HElib and Lattigo.

\end{document}
